\documentclass[11pt]{article}
\usepackage{amsmath,amssymb,booktabs,geometry,array}
\geometry{margin=1in}
\setlength{\parskip}{4pt}
\renewcommand{\arraystretch}{1.3}
\newcommand{\Cframe}{\mathcal{C}}
\newcommand{\Iframe}{\mathcal{I}}
\newcommand{\Lw}{L_\omega}

\title{\textbf{Two-Tier Visibility Conditioning of the Acceleration Command}\\[4pt]
\large A minimal lean projection on the real camera plane, plus a self-releasing descent ease}
\author{Soft--Precise--Landing}
\date{9 September 2026 \quad(supersedes the 31 Aug 2026 joint-QP note)}

\begin{document}
\maketitle

\noindent\emph{Scope.} Keep the marker's \emph{measured centre} inside the \emph{real} camera image
plane, with a fixed buffer from the field-of-view edge, by making the smallest useful change to the
desired inertial acceleration $\,^\Iframe\boldsymbol{a}_\text{d}$ an outer loop produces. For the cross
marker the visibility target is a \emph{single point} --- the crossed-lines intersection --- so no
4-corner containment and no marker-size term. The command reaches that point through two channels on
different timescales, handled as two tiers rather than one fused problem:
\begin{itemize}
\item \textbf{Attitude (lean).} A body tilt rotates the FoV, so the projected centre shifts within one
control step, depth-free. This is the \textbf{hard} constraint (Tier~1).
\item \textbf{Descent rate.} $|\tilde{\boldsymbol{r}}|\approx|\text{offset}|/Z$ grows as $Z$ shrinks, at a
rate $\sim|\tilde{\boldsymbol{r}}|\,(-\dot Z/Z)$. A change in the vertical command only affects the centre
two integrators later, so descent is a \textbf{soft, predictive} influence (Tier~2): if the centre will
reach the buffer before the lateral loop can re-centre it, ease off the descent to buy that loop time.
\end{itemize}
The tilt$\to$feature coupling is the \textbf{rotational part of the IBVS interaction matrix}, $L_\omega$,
evaluated at the measured point --- exact, \emph{depth-independent}, no de-rotation or virtual frame.
The 31~Aug 2026 note posed one QP over $\,^\Iframe\boldsymbol{a}_\text{d}$ jointly enforcing visibility
\emph{and} thrust deliverability with $a_z$ a solved variable; \S\ref{sec:retired} records why that was
retired.

\section*{Notation}
\begin{tabular}{@{}l >{\raggedright\arraybackslash}p{0.66\textwidth}@{}}
\toprule
Symbol & Meaning \\
\midrule
$\tilde{\boldsymbol{r}},\ \tilde r_{k}$ & \emph{measured} camera-frame marker centre, tangent units ($\text{px}/f$), per axis $k$ \\
$\boldsymbol\varphi=\tfrac{\mathcal{R}}{2f}(1-b)$ & FoV-edge tangent half-extent minus the buffer; $b=$ \texttt{CBF\_BUFFER\_FRAC} $=0.15$ \\
$L_\omega(\tilde{\boldsymbol{r}})$ & rotational interaction matrix $\left[\begin{smallmatrix}xy & -(1+x^2)\\ 1+y^2 & -xy\end{smallmatrix}\right]$ at $(x,y)=\tilde{\boldsymbol r}$ --- \emph{no depth} \\
$M=\left[\begin{smallmatrix}0&1\\-1&0\end{smallmatrix}\right]$ & lean$\to$rotation-axis map (fixed $90^\circ$): a lean along $\boldsymbol{y}$ is a rotation about $M\boldsymbol{y}$ \\
$L_e=-(L_\omega M)$ & effective Jacobian $\partial\tilde{\boldsymbol{r}}/\partial\boldsymbol{y}$. \textbf{Sign negated} vs the old $L_\omega M$ (\S\ref{sec:tier1}) \\
$P=R_z(90^\circ)R_z(-\psi_\text{b})$ & inertial $x$-$y$ $\to$ image axes (sim camera mounted yaw$+90^\circ$ on the down-pitch) \\
$\,^\Iframe\boldsymbol{a}_\text{d}$ & desired inertial \emph{specific thrust}, NED. \textbf{Hover $=-g\boldsymbol{e}_3$}; $\|\cdot\|$ \emph{is} thrust/mass. Impl.\ \texttt{I\_a} \\
$a_z=|\,^\Iframe a_{\text{d},z}|$ & vertical magnitude --- a \textbf{fixed input} to Tier~1, never solved \\
$\boldsymbol{y}$ & commanded lean $P(\,^\Iframe\boldsymbol{a}_{\text{d},xy}/a_z)$, image axes; $\boldsymbol{y}_\text{d}$ from $\,^\Iframe\boldsymbol{a}_\text{d}^{\,0}$; $\boldsymbol{y}_\text{curr}=P(-R_{[{:}2,3]}/R_{33})$ from the realized attitude \\
$\tau,\boldsymbol{d}$ & drift look-ahead (s); measured centre drift rate (tangent$/$s, gyro-stripped). $\tau=0$ for a stationary target \\
$A_\text{cap},\ g$ & vehicle max thrust-to-mass ratio; gravity \\
$g_z\in[g_\text{min},1]$ & Tier-2 descent-ease scale (\S\ref{sec:tier2}); $g_\text{min}=$ \texttt{CBF\_GMIN} $=0.2$ \\
$T_\text{react}$ & lateral-loop reaction timescale, \texttt{CBF\_TREACT} $=1.5$\,s \\
\bottomrule
\end{tabular}

\section{Barrier on the real camera plane}
Per image axis $k$,
\begin{equation}\label{eq:barrier}
\boxed{\,h_{k}=\varphi_{k}-|\tilde r_{k}|\,}, \qquad h_k\ge0\iff\text{the centre is on the sensor on axis }k.
\end{equation}
$\tilde{\boldsymbol{r}}$ is the \emph{measured} centre --- nothing is de-rotated. $\boldsymbol\varphi$ is a
fixed camera constant. The buffer $b$ absorbs, in one number, the intended edge margin, the one-cycle
linearization residual of \eqref{eq:N}, and inner-loop attitude-tracking lag.

\section{Tilt$\to$feature coupling, in acceleration units}\label{sec:coupling}
A camera rotation moves the centre by $L_\omega$, a standard \emph{depth-independent} object at the
measured point:
$\dot{\tilde{\boldsymbol{r}}}=L_\omega\boldsymbol\omega_{\rm rp}+\boldsymbol{d}$.
The quantity a lateral-acceleration command sets is not a rotation rate but the underactuated
quadrotor's \emph{lean} $\boldsymbol{y}$; producing a steady lean requires a rotation about the
perpendicular axis, $\boldsymbol\omega_{\rm rp}=M\boldsymbol{y}$. Hence the feature response to the
lean is $L_\omega M\,\boldsymbol{y}$, and to the lateral acceleration directly
\begin{equation}\label{eq:N}
\tilde{\boldsymbol{r}}(\,^\Iframe\boldsymbol{a}_{\text{d},xy})\approx
\tilde{\boldsymbol{r}}+L_e\big(\boldsymbol{y}-\boldsymbol{y}_\text{curr}\big)+\tau\boldsymbol{d},
\qquad \boldsymbol{y}=P\,\frac{\,^\Iframe\boldsymbol{a}_{\text{d},xy}}{a_z},
\end{equation}
with $a_z$ a \emph{given} (the outer loop's commanded value). Equation~\eqref{eq:N} is the local
linearization, accurate over one control step; $\boldsymbol{y}_\text{curr}$ is read from the realized
attitude, so it needs no rotation from acceleration units.

\section{Tier 1 --- the hard, this-cycle lean projection}\label{sec:tier1}
\begin{equation}\label{eq:proj}
\boxed{\;\,^\Iframe\boldsymbol{a}_\text{d}^\star\!:\quad
\boldsymbol{y}^\star=\arg\min_{\boldsymbol{y}}\ \|\boldsymbol{y}-\boldsymbol{y}_\text{d}\|^2
\quad\text{s.t.}\quad
\big|\tilde{\boldsymbol{r}}+L_e(\boldsymbol{y}-\boldsymbol{y}_\text{curr})+\tau\boldsymbol{d}\big|_k\le\varphi_k
\ \ (k=1,2)\;}
\end{equation}
solved by a few alternating projections of $\boldsymbol{y}_\text{d}$ onto the two affine rows, then
$\,^\Iframe\boldsymbol{a}_{\text{d},xy}^\star=a_z\,P^{\!\top}\boldsymbol{y}^\star$ and
$\,^\Iframe a_{\text{d},z}^\star=\,^\Iframe a_{\text{d},z}$ \textbf{unchanged}.

\begin{itemize}
\item \textbf{Minimal intervention.} A row is touched only when it is \emph{predicted} to breach. If
$\boldsymbol{y}_\text{d}$ already keeps the centre in frame, $\boldsymbol{y}^\star=\boldsymbol{y}_\text{d}$
exactly.
\item \textbf{Outward-only.} Leaning \emph{toward} the marker (or tangentially) reduces $|\tilde r_k|$, so
no row fires --- a re-centring command is never clipped.
\item \textbf{Depth-free, scale-free.} Every term is a camera intrinsic, an attitude, an image-space
measurement, or a ratio/pass-through of the caller's own command.
\end{itemize}

\paragraph{Jacobian sign.} $L_e=-(L_\omega M)$. The extra minus is not in the textbook
forward-camera form; it was fixed against an \emph{independent} pinhole + attitude camera model
(\texttt{tools/validate\_visibility\_projection.py} check~1): for the \emph{downward} camera the feature
moves opposite to $L_\omega M\,\boldsymbol{y}$. The retired machinery used $L_\omega M$; its own offline
validator carried a matching historical sign error, so the mistake was self-consistent and unseen.

\paragraph{Linearization envelope.} \eqref{eq:N} uses $L_\omega$ at the \emph{measured} centre, fixed
for the solve. It is $<8\%$ off for a per-cycle lean \emph{change} up to $\sim15^\circ$ --- the real
regime at 50\,Hz with a smooth outer loop --- and the buffer $b$ absorbs the residual (offline: 0/3742
true-sensor exits at $\le15^\circ$/cycle; 13/3762 at $\sim24^\circ$, worst $+0.07$ tangent). It is
\emph{not} accurate for a one-shot $\sim50^\circ$ correction; the every-cycle incremental nature of the
correction and the lean cap below cover that.

\paragraph{Deliverability is separate.} Thrust deliverability is \emph{not} in \eqref{eq:proj}. The
caller applies it afterward: a lean cap $\|\boldsymbol{y}^\star\|\le\arccos(a_z/A_\text{cap})$ (the true
deliverable tilt at the actual $a_z$; $\theta_\text{cap}$ if $A_\text{cap}\le g$), then a thrust-vector
cap $\|\,^\Iframe\boldsymbol{a}_\text{d}\|\le A_\text{cap}$. Folding these into the projection is what
forced $a_z$ to become a solved variable and created the coupled iteration of the old note; keeping them
out keeps $a_z$ a fixed input and the solve a plain feasibility projection.

\section{Tier 2 --- the soft, self-releasing descent ease}\label{sec:tier2}
While the box binds it suppresses lateral authority the outer loop asked for. Pressing on at the full
closing rate spends the descent budget before the lateral error is corrected. Tier~2 spends part of that
budget on time, but keyed on a \emph{measured} quantity, not on how much the box clipped:
\begin{equation}\label{eq:ease}
\text{rem}_k=\varphi_k-|\tilde r_k|,\quad
\text{close}_k=\max\!\big(\operatorname{sgn}(\tilde r_k)\,\dot{\tilde r}_k,\,0\big),\quad
t_\text{edge}=\min_k\frac{\text{rem}_k}{\text{close}_k},
\end{equation}
\begin{equation}\label{eq:gz}
\boxed{\;g_z=g_\text{min}+(1-g_\text{min})\operatorname{clip}\!\Big(\frac{t_\text{edge}}{T_\text{react}},0,1\Big),
\qquad
\,^\Iframe a_{\text{d},z}^{\star\star}=(\,^\Iframe a_{\text{d},z}^\star+g)\,g_z-g\;}
\end{equation}
applied only while descending ($\,^\Iframe a_{\text{d},z}^\star+g>0$) and low-passed so $a_z$ does not
step. $t_\text{edge}$ is an image-distance over an image-speed, i.e.\ \emph{seconds} --- scale-free.

\paragraph{Why this is safe where the old descent relief was not.} The retired
\texttt{CBF\_AZ\_COST\_GAIN} relief triggered on $\|\boldsymbol{y}_\text{d}-\boldsymbol{y}^\star\|$, a
QP-internal norm that self-inflated as the relief's own prior-iterate lift changed the $a_z$ it was
measured against, and it acted through $\min(\cdot,-g)$ which could \emph{pin} $a_z$ at hover
(descent-freeze deadlock). Tier~2 instead: (i)~triggers on an external measured signal ($|\tilde r|$ and
its rate); (ii)~runs \emph{after} Tier~1, so it never fights the lean solve; (iii)~is a \emph{bounded
scale} with $g_\text{min}>0$ --- it can only slow a descent, and always leaves some; (iv)~\textbf{self-
releases}: the bought time shrinks $|\tilde r|$, so $\text{rem}$ grows, $t_\text{edge}$ grows, $g_z\to1$.

\section{Feasibility}
\eqref{eq:proj} is infeasible iff $\varphi_k<0$ for some $k$ (the marker subtends the full FoV minus the
buffer). Otherwise the alternating projection always converges. The large-tilt case where the linearized
constraint and the true rotation homography diverge is not handled by a feasibility test but by the
buffer plus the fact that the correction runs every 50\,Hz cycle on a smooth command --- a single cycle
never demands a large lean change.

\section{Assumptions}
\begin{enumerate}
\item The coupling $L_e=-(L_\omega M)$ and its sign are calibrated (validator check~1 vs an independent
pinhole model).
\item The conditioned command is realized: the inner attitude loop tracks the attitude built from
$\boldsymbol{y}^\star$ on a timescale fast relative to the outer loop; residual lag is inside the buffer.
\item $\boldsymbol{d}$ is the translation part of the observed centre velocity, gyro-stripped; $\tau=0$
for a stationary target (moving-target wiring is a separate step).
\item First-order $L_\omega$ over the cycle (envelope $\sim15^\circ$/cycle; \S\ref{sec:tier1}).
\item $A_\text{cap},g$ are fixed vehicle/physical constants, used directly.
\end{enumerate}

\section{Scale-freeness and depth-freeness}
\eqref{eq:barrier}, \eqref{eq:N}, \eqref{eq:proj}, \eqref{eq:ease}--\eqref{eq:gz} contain only the
measured image point, $\boldsymbol\varphi=\mathcal{R}/(2f)\,(1-b)$, the body rate, the commanded
acceleration and its measured image-space rate, $A_\text{cap}$, $g$, and fixed time constants. $L_\omega$
is depth-independent by construction; $t_\text{edge}=\text{(image distance)}/\text{(image speed)}$ is
dimensionless in scale. The formulation is manifestly scale- and depth-free.

\section{What this retired, and why}\label{sec:retired}
The 31~Aug 2026 joint QP over $\,^\Iframe\boldsymbol{a}_\text{d}$ (visibility box + deliverability sphere,
$a_z$ solved), its two-phase $\delta$ / decode-fail geometry QP, the descent-rate relief
(\texttt{CBF\_AZ\_COST\_GAIN}), the frozen $\rho_\text{fov}$ / $d_\text{min,fov}$ tilt-cone and its
overflow/drift-off classification, and the ArUco 4-corner copy were all removed
(\texttt{cbf\_visibility.py}, \texttt{cbf\_visibility\_aruco.py} $\to$ \texttt{Obsolete/}).
Grounds:
\begin{itemize}
\item The deliverability \textbf{sphere never bound} in flight (0.0\% of gate frames; 81--97\% of
$A_\text{cap}$) --- the post-step $\arccos(a_z/A_\text{cap})$ lean cap already covers it.
\item With the sphere idle and the relief off, the joint projection is \textbf{algebraically the plain
tilt projection} ($N_i(a_z)\,^\Iframe\boldsymbol{a}_{\text{d},xy}\equiv L_\omega M\,\boldsymbol{y}$) ---
the $a_z$-in-the-loop iteration bought nothing.
\item The descent relief \textbf{failed every SITL gate} in three forms; a clean IC2--5 A/B put
\texttt{CBF\_AZ\_COST\_GAIN}$=0$ ahead of $=5$ (tighter, softer, no regression).
\item The $\rho_\text{fov}$ cone was \textbf{deliberately frozen} (\texttt{PLASMC\_LFOV}$=0$) and dormant.
\item The two-phase $\delta$ apparatus fired \textbf{0.1\% of frames} on the live cross-marker pipeline.
\end{itemize}
The only retained consumer of the old classification, the drift-off pull-back, is re-expressed as a
per-axis increase of $b$ on a persistent one-sided $\varphi$ breach.

\paragraph{Validation of this design.} \texttt{tools/validate\_visibility\_projection.py} --- 10/10
against an independent pinhole + attitude oracle (barrier holds, minimal intervention, outward-free,
conventions round-trip, descent-ease one-way / triggered-by-closing / self-releasing). SITL: IC2--5 $n=5$
cross-marker A/B vs the retired machinery --- both 20/20 land, 0 target-loss; this design pooled mean
$xy$ $0.14$\,m vs $0.24$, precise $11/20$ vs $8/20$, max touchdown speed $1.47$ vs $1.95$\,m/s. The
median is a wash (on a clean approach the projection is a pure pass-through); the gain is entirely in the
tail, where the retired joint-solve / relief / pull-back / cone-floor stack was found to thrash
(intervening against the outer loop's re-centring, chattering the tilt into a gain ratchet) on exactly
the marginal off-centre approaches the gate stresses.

\end{document}
